\section{Proof of the Delay Theorem}\label{app:proof}

\begin{proof}[Proof of \cref{thm:delay}]
We compare the real experiment with an ideal one in which the values $\Open$ computes are drawn only when the attacker reaches them, show that the ideal experiment is sequential, and bound the probability that the two differ.

\emph{Path values.}
For each chain $i$, let $y_{i,0} = s_i = \Hstart(\ctx, i, x_i)$ and $y_{i,k} = F_i(y_{i,k-1})$ for $k = 1, \dots, m$, so that $y_{i,m} = y_i$. Let $\mu_i = \Hmask(\ctx, y_i)$ for $i < n$, $a_0 = \ctx$, $a_i = \Hacc(a_{i-1}, y_i)$, $K = \Hkey(\ctx, a_n)$ and $C = \Hcommit(\ctx, K)$. These $P$ values are the \emph{path values}. Each path value $v$ is the answer on one \emph{path input} $u(v)$. The \emph{predecessors} of $v$ are the path values in $u(v)$, at most two; in addition, $\mu_{i-1}$ is a predecessor of $s_i$ for $i \ge 2$, since $x_i = z_{i-1} \oplus \trunc{\mu_{i-1}}{b}$. Let $\rho(v)$ be the oracle round in which $\Open$ receives $v$:
\[
  \rho(s_i) = (i - 1)(m + 2) + 1, \qquad
  \rho(y_{i,k}) = \rho(s_i) + k, \qquad
  \rho(\mu_i) = \rho(a_i) = \rho(y_i) + 1, \qquad
  \rho(K) = \Tkey.
\]
Every path value $v$ other than $s_1$ satisfies $\rho(v) = 1 + \max_w \rho(w)$, the maximum taken over its predecessors $w$.

\emph{The two experiments.}
In the real experiment $\Real$, the oracles are random functions sampled lazily; $\Adv_0$ runs, $\Build$ runs, and $\Adv_1$ receives $Z$. The ideal experiment $\Ideal$ uses the same lazily sampled oracles, but $\Build$ makes no queries. It draws the salt and $x_1$ as before, and draws $z_1, \dots, z_{n-1}$ and $C$ uniformly. A query of $\Adv_1$ is \emph{legitimate} if, when it is submitted, it equals $u(v)$ for a path value $v$ whose predecessors have all been revealed; for $s_i$ with $i \ge 2$ this defines $x_i = z_{i-1} \oplus \trunc{\mu_{i-1}}{b}$. $\Ideal$ answers a legitimate query with $v$, drawing $v$ uniformly the first time ($C$ was drawn already), and marks $v$ revealed. It answers every other query, and all queries of $\Adv_0$, from the lazily sampled oracles. When $\Adv_1$ stops, $\Ideal$ draws the path values never revealed, which fixes the remaining links.

\begin{claim}\label{claim:sequential}
In $\Ideal$, no path value $v$ is revealed before oracle round $\rho(v)$.
\end{claim}
\begin{claimproof}
By induction on $\rho(v)$; for $s_1$, with $\rho = 1$, there is nothing to show. A value is revealed only by a legitimate query, which is submitted after all its predecessors are revealed. Answers arrive at the end of a round, so if a predecessor $w$ is revealed in round $\rho(w)$ or later, $v$ is revealed in round $\rho(w) + 1$ or later; hence in round $1 + \max_w \rho(w) = \rho(v)$ or later.
\end{claimproof}

\begin{claim}\label{claim:hybrid}
Let $\Bad$ be the event, in $\Ideal$, that (a)~two path values are equal, or (b)~a query that is not legitimate equals a path input, evaluated with the path values as finally drawn. Then the probability that $\Adv_1$ outputs $K$ early in $\Real$ is at most the probability that it does so in $\Ideal$ plus $\Pr_{\Ideal}[\Bad]$.
\end{claim}
\begin{claimproof}
Run both experiments on the same randomness, with the oracle's answers on path inputs in $\Real$ taken to be the path values drawn in $\Ideal$. If (a) fails, the path inputs are distinct, so these answers are independent and uniform as in $\Real$; the puzzle has the same distribution in both, because $z_i = \trunc{\mu_i}{b} \oplus x_{i+1}$ is uniform when $x_{i+1}$ is, and $C$ is an answer on a fresh input. If (b) fails, every query either is legitimate, and receives the path value in both experiments, or is not a path input, and receives the same oracle answer in both. So the two transcripts agree unless $\Bad$ happens.
\end{claimproof}

\begin{claim}\label{claim:bad}
The probability of\/ $\Bad$ in\/ $\Ideal$ is at most
\[
  \Pr\nolimits_{\Ideal}[\Bad] \le \frac{\Qon}{2^{8b}} + \frac{\Qpre}{2^{128 + 8b}} + \frac{QP + P^2}{2^{256}}.
\]
\end{claim}
\begin{claimproof}
In $\Ideal$, a path value that is not revealed when a query is submitted is drawn later, independently of that query. So is $x_i$ for $i \ge 2$ while $\mu_{i-1}$ is not revealed, and so are the salt and $x_1$ for every query of $\Adv_0$. Nothing is conditioned on: each bound below is the probability that a fixed query equals a value drawn after it.
\begin{itemize}[leftmargin=1.5em]
  \item Event (a): there are at most $P^2$ pairs of independent uniform 32-byte values, so its probability is at most $P^2 / 2^{256}$.
  \item Event (b), queries to $\Hstart$: a query names one index $i$ and equals $u(s_i)$ only if it contains $x_i$. A query of $\Adv_1$ with $i = 1$ is legitimate whenever it equals $u(s_1)$. One with $i \ge 2$ that is not legitimate was submitted while $\mu_{i-1}$ was not revealed, so it contains $x_i$ with probability $2^{-8b}$; this case needs $n \ge 2$. A query of $\Adv_0$ must also contain the salt, so it succeeds with probability $2^{-(128 + 8b)}$.
  \item Event (b), all other queries: a query that is not legitimate but equals $u(v)$ contains, in a 32-byte field, a predecessor of $v$ that was not revealed when it was submitted. That happens with probability at most $2^{-256}$ for each of the at most $P$ path inputs, so at most $P / 2^{256}$ per query.
\end{itemize}
Summing over the $\Qon$ and $\Qpre$ queries gives the claim.
\end{claimproof}

\emph{Conclusion.}
In $\Ideal$, if $\Adv_1$ stops before round $\Tkey$, then $K$ is not revealed by \cref{claim:sequential}, so it is drawn after $\Adv_1$'s output, and the output equals it with probability $2^{-256}$; $C$, drawn independently, tells $\Adv$ nothing about it. \Cref{claim:hybrid,claim:bad} add $\Pr_{\Ideal}[\Bad]$.
\end{proof}

The proof follows the structure of Pietrzak and Reyzin's soundness proof for proofs of sequential work (their Lemmas 6 and~9)~\cite{PR26}. Section~8 of the specification gives the same argument for \leanseal's byte encoding. It has not been machine-checked or reviewed by a third party.
